\subsection{Dimension of a quotient ring}

PROPOSITION 1. --- Let A be a Noetherian integral domain,  a nonzero element of A, and  a minimal element of the set of prime ideals of A containing . Then  has height 1.

Let \(\mathfrak{q} \subset \mathfrak{p}\) be a prime ideal distinct from \(\mathfrak{p}\). We have \(x \notin \mathfrak{q}\) by the minimality of \(\mathfrak{p}\). Since A is an integral domain, \(\mathrm{A}_{\mathfrak{p}}\) is identified with a subring of \(\mathrm{A}_{\mathfrak{q}}\); for every integer \(n \geqslant 0\), denote by \(\mathfrak{q}_n\) the ideal \(\mathfrak{q}^n\mathrm{A}_{\mathfrak{q}} \cap \mathrm{A}_{\mathfrak{p}}\) of \(\mathrm{A}_{\mathfrak{p}}\). The minimality of \(\mathfrak{p}\) means that the local ring \(\mathrm{A}_{\mathfrak{p}} / x\mathrm{A}_{\mathfrak{p}}\) has dimension 0; it is therefore of finite length (§ 1, no. 3, example 1), and there exists an integer \(n_0 \geqslant 0\) such that one has

\[
\mathfrak {q} _ {n} + x \mathrm{A} _ {\mathfrak {p}} = \mathfrak {q} _ {n + 1} + x \mathrm{A} _ {\mathfrak {p}} \quad \text { for   tout } \quad n \geqslant n _ {0}.\tag{1}
\]

Fix the integer \(n \geqslant n_{0}\) . Given \(y \in \mathfrak{q}_{n}\) , there exists \(a \in \mathrm{A}_{\mathfrak{p}}\) such that \(y - ax \in \mathfrak{q}_{n+1}\) ; we then have \(ax \in \mathfrak{q}_{n}\) , whence \(a \in \mathfrak{q}_{n}\) since \(x \notin \mathfrak{q}\) , and finally we have \(y \in \mathfrak{q}_{n+1} + x\mathfrak{q}_{n}\) . One therefore has

\[
\mathfrak {q} _ {n} = \mathfrak {q} _ {n + 1} + x \mathfrak {q} _ {n}.\tag{2}
\]

Since  belongs to the maximal ideal of the Noetherian local ring , Nakayama's lemma shows that one has \(A_{p}\)\(\mathfrak{q}_n = \mathfrak{q}_{n+1}\). Since one has \((\mathfrak{q}A_{\mathfrak{q}})^{n} = \mathfrak{q}_{n}A_{\mathfrak{q}}\), one concludes

\[
(q A _ {q}) ^ {n} = (q A _ {q}) ^ {n + 1} \quad \text { for   tout } \quad n \geqslant n _ {0}.\tag{3}
\]

Since the ring \(\mathbf{A}_{\mathfrak{q}}\) is local and Noetherian, one has \(\bigcap_{n\geqslant 0}(\mathfrak{q}\mathbf{A}_{\mathfrak{q}})^n = \{0\}\) (III, § 3, no. 2, corollary of prop. 5), whence \((\mathfrak{q}\mathbf{A}_{\mathfrak{q}})^{n_0} = \{0\}\) and finally the prime ideal \(\mathfrak{q}\mathbf{A}_{\mathfrak{q}}\) of \(\mathbf{A}_{\mathfrak{q}}\) is reduced to 0. We therefore have \(\mathfrak{q} = \{0\}\) , which proves that \(\mathfrak{p}\) is of height 1.

PROPOSITION 2. --- Let A be a Noetherian ring,  a positive integer, and  an ideal contained in the radical of A and generated by  elements. Then one has

\[
\dim (\mathrm{A} / \mathfrak {a}) \leqslant \dim (\mathrm{A}) \leqslant \dim (\mathrm{A} / \mathfrak {a}) + m.\tag{4}
\]

The inequality  follows from Proposition 6 of § 1, no. 3. An immediate induction on  shows that it suffices to establish the inequality

\[
\dim (\mathrm{A}) \leqslant \dim (\mathrm{A} / x \mathrm{A}) + 1\tag{5}
\]

for every element \(x\) of the radical of A, that is, to prove that one has \(\dim(A/xA) \geqslant n - 1\) for every chain \(\mathfrak{p}_0 \subset \ldots \subset \mathfrak{p}_n\) of prime ideals of A, of length \(n \geqslant 1\), and such that \(x \in \mathfrak{p}_n\). It suffices to construct a chain \(\mathfrak{q}_1 \subset \ldots \subset \mathfrak{q}_n\) of prime ideals of A, with \(x \in \mathfrak{q}_1\), and this follows from the following lemma:

Lemma 1. --- Let A be a Noetherian ring, \(\mathfrak{p}_{0} \subset \ldots \subset \mathfrak{p}_{n}\) a chain of prime ideals of A of length \(n \geqslant 1\), and \(x\) an element of \(\mathfrak{p}_{n}\). There exists a chain \(\mathfrak{p}_{0}' \subset \ldots \subset \mathfrak{p}_{n}'\) with \(\mathfrak{p}_{0}' = \mathfrak{p}_{0}\), \(\mathfrak{p}_{n}' = \mathfrak{p}_{n}\), and \(x \in \mathfrak{p}_{1}'\).

Let us reason by induction on \(n\), the case \(n=1\) being trivial. Suppose therefore that we have \(n\geqslant2\) and that \(x\) does not belong to \(\mathfrak{p}_{n-1}\). Let \(\mathfrak{p}_{n-1}^{\prime}\) be a minimal element of the set of prime ideals of A contained in \(\mathfrak{p}_{n}^{\prime}=\mathfrak{p}_{n}\) and containing \(\mathfrak{p}_{n-2}+\mathbf{A}x\) (II, § 2, no. 6, lemma 2). By Proposition 1, the ideal \(\mathfrak{p}_{n-1}^{\prime}/\mathfrak{p}_{n-2}\) of the ring \(\mathbf{A}/\mathfrak{p}_{n-2}\) has height 1, and since \(\mathfrak{p}_{n-2}\subset\mathfrak{p}_{n-1}\subset\mathfrak{p}_{n}\) is a chain of length 2, the same is true of \(\mathfrak{p}_{n-2}\subset\mathfrak{p}_{n-1}^{\prime}\subset\mathfrak{p}_{n}^{\prime}\). We have \(x\in\mathfrak{p}_{n-1}^{\prime}\). The induction hypothesis applied to the chain \(\mathfrak{p}_{0}\subset\mathfrak{p}_{1}\subset\ldots\subset\mathfrak{p}_{n-2}\subset\mathfrak{p}_{n-1}^{\prime}\) shows that there exists a chain \(\mathfrak{p}_{0}^{\prime}\subset\mathfrak{p}_{1}^{\prime}\subset\ldots\subset\mathfrak{p}_{n-2}^{\prime}\subset\mathfrak{p}_{n-1}^{\prime}\) with \(x\in\mathfrak{p}_{1}^{\prime}\) and \(\mathfrak{p}_{0}^{\prime}=\mathfrak{p}_{0}\). The chain

\[
\mathfrak {p} _ {0} ^ {\prime} \subset \mathfrak {p} _ {1} ^ {\prime} \subset \dots \subset \mathfrak {p} _ {n - 1} ^ {\prime} \subset \mathfrak {p} _ {n} ^ {\prime}
\]

satisfies the required conditions.

COROLLARY 1. --- a) Every Noetherian local ring has finite dimension. More generally, every nonzero Noetherian semilocal ring (II, § 3, no. 5, definition 4) has finite dimension.

\begin{enumerate}
\def\labelenumi{\alph{enumi})}
\setcounter{enumi}{1}
\item
  Let A be a Noetherian ring. Every ideal of A other than A has finite height.
\item
  Every decreasing sequence of prime ideals of a Noetherian ring A is stationary.
\item
  Let A be a nonzero Noetherian semilocal ring and let  be its radical; the quotient ring  is Artinian and nonzero, hence of dimension 0 (§ 1, no. 3, example 1). There exists an integer  such that the ideal  of A is generated by  elements; we therefore have  by Proposition 2.
\item
  Let  be an ideal of A, and let \(\mathfrak{a} \neq A\) be a maximal ideal of A containing \(\mathfrak{a}\). We have  by Proposition 7 of § 1, no. 3, and \(0 \leqslant \mathrm{ht}(\mathfrak{a}) \leqslant \dim(A_{\mathfrak{m}})\)\(A_{\mathfrak{m}}\) is a Noetherian local ring. Therefore \(\mathrm{ht}(\mathfrak{a})\) is finite by a).
\item
  Every finite strictly decreasing sequence \((\mathfrak{p}_{i})_{0\leqslant i\leqslant n}\) of prime ideals of A defines a chain , whence \(\mathfrak p_{n}\subset\ldots\subset\mathfrak p_{0}\)\(n\leqslant\dim(A_{\mathfrak p_{0}})<+\infty\). There can therefore not exist an infinite strictly decreasing sequence of prime ideals of A, whence c).
\end{enumerate}

COROLLARY 2. --- Let A be a Noetherian local ring.

\begin{enumerate}
\def\labelenumi{\alph{enumi})}
\item
  Let \(x \in \mathfrak{m}_{\mathbf{A}}\). Then \(\dim(\mathbf{A}/x\mathbf{A})\) is equal to \(\dim(\mathbf{A})\) or to \(\dim(\mathbf{A}) - 1\). For one to have \(\dim(\mathbf{A}/x\mathbf{A}) = \dim(\mathbf{A}) - 1\), it is necessary and sufficient that \(x\) does not belong to any of the minimal prime ideals \(\mathfrak{p}\) of \(\mathbf{A}\) such that \(\dim(\mathbf{A}/\mathfrak{p}) = \dim(\mathbf{A})\), and it suffices that \(x\) is not a zero divisor in \(\mathbf{A}\).
\item
  Let \(\mathfrak{a}\) be an ideal of A other than A such that \(\dim(\mathbf{A}/\mathfrak{a}) < \dim(\mathbf{A})\). There exists \(x \in \mathfrak{a}\) such that \(\dim(\mathbf{A}/x\mathbf{A}) = \dim(\mathbf{A}) - 1\).
\item
  If , there exists  such that \(	ext{If} \dim(A) \geqslant 1\)\(there exists x \in \mathfrak{m}_{A} \text{ \text{such that} } \dim(A/xA) = \dim(A) - 1\).
\end{enumerate}

By Proposition 2,  is equal to  or to . For \(\dim(A/xA) = \dim(A)\) to hold, it is necessary and sufficient that there exist a chain \(\mathfrak{p}_0 \subset \ldots \subset \mathfrak{p}_n\) of prime ideals of A such that \(x \in \mathfrak{p}_0\) and \(n = \dim(A)\), that is, there exists a prime ideal \(\mathfrak{p}_0\) of A containing \(x\) such that \(\dim(A/\mathfrak{p}_0) = \dim(A)\). But such a prime ideal \(\mathfrak{p}_0\) is necessarily minimal, and every element of \(\mathfrak{p}_0\) is therefore a zero divisor in A (IV, § 1, no. 1, cor. 3 of Proposition 2 and no. 4, theorem 2). This proves \(a)\).

Let \(\Phi\) be the set of minimal prime ideals of A, and \(\Phi'\) the set of \(\mathfrak{p} \in \Phi\) such that \(\dim(A/\mathfrak{p}) = \dim(A)\). We know (II, § 4, no. 3, cor. 3 of Proposition 14) that \(\Phi\) is finite, hence \(\Phi'\) is finite. Let \(\mathfrak{a}\) be an ideal of A such that \(\dim(A/\mathfrak{a}) < \dim(A)\). For every \(\mathfrak{p} \in \Phi'\), one has \(\dim(A/\mathfrak{a}) < \dim(A/\mathfrak{p})\), hence \(\mathfrak{a} \not\subsetneq \mathfrak{p}\). By Proposition 2 of II, § 1, no. 1, there therefore exists an element \(x\) of \(\mathfrak{a}\) which belongs to none of the \(\mathfrak{p} \in \Phi'\), and one then has \(\dim(A/xA) = \dim(A) - 1\) by \(a)\). This proves \(b)\).

Assertion c) is the particular case \(\mathfrak{a} = \mathfrak{m}_{\mathrm{A}}\) of \(b\) ).
% semantic-fragments: legacy-input-seam
\relax{}
